\pdfoutput=1
\documentclass[11pt]{article}
\usepackage[margin=1in]{geometry}
\usepackage[T1]{fontenc}
\usepackage{amsmath,amssymb,amsthm,booktabs,tabularx}
\usepackage{algorithm}
\usepackage[noend]{algpseudocode}
\usepackage[hidelinks]{hyperref}

\newtheorem{theorem}{Theorem}
\newtheorem{lemma}[theorem]{Lemma}
\newtheorem{corollary}[theorem]{Corollary}
\theoremstyle{remark}
\newtheorem{remark}[theorem]{Remark}

\newcommand{\lean}[1]{\texttt{\detokenize{#1}}}
\newcommand{\SIZED}{\textsf{SIZED}}
\newcommand{\COUNT}{\textsf{COUNT}}
\newcommand{\DOMCOUNT}{\textsf{DOMCOUNT}}
\newcommand{\poly}{\operatorname{poly}}
\newcommand{\cnt}{\mathrm{cnt}}
\renewcommand{\algorithmicrequire}{\textbf{Input:}}
\renewcommand{\algorithmicensure}{\textbf{Output:}}
\algrenewcommand\algorithmicindent{1.2em}

\title{Machine-Checked Algorithmic Consequences\\ of a \texorpdfstring{$2^{o(n)}$}{2\textasciicircum o(n)} Algorithm for \texorpdfstring{\#$k$}{\#k}-SAT}
\hypersetup{pdftitle={Machine-Checked Algorithmic Consequences of a 2\textasciicircum o(n) Algorithm for \#k-SAT},pdfauthor={Sebastian Senge}}
\author{Sebastian Senge\\ Independent Researcher}
\date{}

\begin{document}
\maketitle

\begin{abstract}
In October 2026 two groups claimed algorithms for \#$k$-SAT in time $2^{o(n)}$, which would refute
the Exponential Time Hypothesis, and both formalized their algorithms in Lean. A refuted lower bound
yields no algorithm by itself: every consequence needs a reduction to \#SAT whose number of
variables is linear in the relevant parameter. We give such reductions for twenty-one problems and
prove each resulting algorithm as a Lean~4 theorem about a single Turing machine, on top of
Anthropic's formalization and with the standard axioms only. The new results are: the parity of
the number of connected vertex sets of size $\ell$ that satisfy a given CNF, on graphs of any
density; from it, with an isolation lemma for arbitrary sets, one-sided-error algorithms for
Steiner tree and connected vertex cover on graphs of any density; the parity and a one-sided-error
decision for connected dominating set on graphs of any density, whose clauses of domination have
unbounded width, by running the clause-shortening tree of Anthropic's dominating-set algorithm
inside the Cut \& Count formula, and from it Max Leaf Spanning Tree and minimum connected dominating
set; the same for any positive CNF of unbounded width with linearly many clauses, which covers Group
Steiner tree, connected hitting set and connected total domination, and the minimum versions of
Steiner tree, connected vertex cover and Group Steiner tree; \#CNF-SAT of unbounded clause
width with up to $n\cdot 2^{(\lg n)^{1/4}/(\lg\lg n)^2}$ clauses; counting $k$-colorings for all
$k$ with $\lg k\le(\lg n)^{1/6-\theta}$; and the Max-SAT spectrum. We also formalize direct
corollaries (homomorphisms into a target in the input, the density of states of the Potts model and
the Tutte polynomial on the hyperbolas $(x-1)(y-1)=q$, the permanent of $\{-1,0,1\}$ matrices in
$2^{\varepsilon m}\poly(n)$) and known consequences (\#CNF-SAT of any width with linearly many
clauses, the 0/1 permanent, Vertex Cover and Nonblocker in $2^{\varepsilon k}\poly(n)$). Our results
on Hamiltonian cycles, the parity for $o(n\sqrt{\lg n}/\lg\lg n)$ edges and one-sided error for
$O(n)$ edges, were obtained independently of Tulsyan's concurrent exact count, which supersedes
them. All bounds hold for every $\varepsilon>0$. We also give reference implementations of all
reductions, tested against brute force and against the Lean definitions.

\end{abstract}

%======================================================================
\section{Introduction}

The Exponential Time Hypothesis (ETH) of Impagliazzo and Paturi~\cite{IP01} states that 3-SAT on
$n$ variables needs time $2^{\Omega(n)}$. Much of fine-grained complexity rests on it: tight lower
bounds for coloring, Hamiltonian cycle and Max-Cut, the $2^{O(\sqrt n)}$ barrier for planar
problems, and most lower bounds for parameterized problems. In October 2026 two groups claimed that
ETH is false.

\begin{itemize}
\item OpenAI~\cite{openai} claims \#$k$-SAT for every fixed $k$ in time
\[(n+m)^{O_k(1)}\cdot 2^{O_k(n\lg\lg n/\sqrt{\lg n})},\]
deterministic and in polynomial space, by contracting Boolean tensor networks with $s$ wires in
time $2^{O(s/\sqrt{\log s})}$.
\item Anthropic~\cite{anthropic} claims \#$k$-SAT with one machine for all $k$ in time
\[2^{C\min(k^2 n\lg\lg n/\sqrt{\lg n},\; km/\sqrt{\lg m})}\]
for an absolute constant $C$, by sum-product computations on the clause graph.
\end{itemize}

Both papers are unrefereed, and both come with Lean formalizations~\cite{formalmath}. If the claims
hold, ETH, SETH, \#ETH, Gap-ETH and NSETH fall. The OV conjecture, the CNF-SAT hypothesis for
$\poly(n)$ clauses of unbounded width, the $k$-Clique conjecture and the Set Cover conjecture still
stand. The algorithms are galactic: Anthropic reports no gain over $2^n$ below $n\approx 10^{10}$.

\paragraph{Corollaries are not automatic.} A refuted lower bound gives no algorithm. To obtain a
$2^{o(n)}$ algorithm for a problem $X$ from a $2^{o(n)}$ algorithm for \#SAT, one needs a reduction
from $X$ to \#SAT with $o(n\sqrt{\lg n}/\lg\lg n)$ variables and, for a uniform bound, small clause
width. Problems in size-constrained SNP reduce in this way~\cite{IPZ01}. Problems with global
constraints (connectivity, orderings, large numbers) or dense pairwise interactions do not, at least
not directly. This paper is about the second kind of problem: we find reductions that keep the
number of variables linear in the parameter even though the natural encoding does not.

\paragraph{Contributions.}
\begin{enumerate}
\item \emph{Connectivity on graphs of any density} (Section~\ref{sec:conn}). The parity of the
number of connected vertex sets of size $\ell$ that satisfy a CNF (N2), and from it one-sided-error
algorithms for Steiner tree and connected vertex cover (N2$'$), on graphs with any number of edges.
The same for connected dominating sets (N7, N7$'$), and from it Max Leaf Spanning Tree and minimum
connected dominating set (N7$''$). More generally, the same for the connected solutions of any positive
CNF of unbounded width with linearly many clauses (N8, N8$'$): Group Steiner tree, connected hitting
set and connected total domination, and the minimum versions (N8$''$).
These are, to our knowledge, the first consequences of the claimed refutations for connectivity
problems on dense graphs; the concurrent work on Hamiltonian cycles~\cite{Tul26} needs average degree
$2^{o(\sqrt{\lg n})}$.
\item \emph{Counting beyond the bounded-width case} (Section~\ref{sec:counting}): \#CNF-SAT of
unbounded width with superlinearly many clauses (A1$'$), $k$-colorings for $k$ growing with $n$ (A2)
and the Max-SAT spectrum (A4), with direct corollaries for homomorphisms, the Potts model and the
Tutte polynomial, and permanents.
\item \emph{Seventeen algorithms} (Table~\ref{tab:results}), each one a polynomial-time reduction to
one of three counters from Anthropic's development, followed by a simple post-step. They are written
out as pseudocode in Sections~\ref{sec:conn}--\ref{sec:fpt}.
\item \emph{Machine-checked statements.} Each algorithm is a Lean~4 theorem stating that one Turing
machine (Mathlib's multi-stack model TM2, finite alphabets) halts on a given input with a given
output within a given number of steps. The proofs build on Anthropic's development, which proves its
$2^{o(n)}$ algorithm as a theorem; the results are therefore Lean theorems, not conditional
statements. They use only the axioms \lean{propext}, \lean{Classical.choice} and \lean{Quot.sound}.
\item \emph{The isolation lemma of~\cite{CIKP08} for arbitrary sets}, formalized: sparse parity
checks on $k$ coordinates isolate a point of any nonempty $S\subseteq\{0,1\}^M$ with probability
$2^{-\delta M}$. The proof of~\cite{CIKP08} is oblivious to the formula~\cite[\S9.3]{Calabro09};
Tulsyan~\cite{Tul26} formalizes the same lemma. Unlike Valiant--Vazirani hashing~\cite{VV86}, it
keeps the number of variables linear.
\item \emph{Reference implementations} of all reductions in Python. They build the same formulas
from the same input words as the Lean machines, and they are tested against brute force and against
the Lean definitions.
\end{enumerate}

\paragraph{Techniques.} Three ideas recur. First, Anthropic's \emph{sized} counter, which counts
the solutions with exactly $s$ true variables, can be combined with complement bits $q=\neg p$: every
complementary pair contributes exactly one true variable, so auxiliary variables leave the target
size fixed. This is how side bits, isolation checks and parity chains enter a sized count
(N2, N2$'$, N6$'$). Second, Cut \& Count~\cite{CNPPRW11} turns connectivity into the parity of a
count, and with a sized counter the number of edges no longer matters. Third, search trees (clause
shortening for A1 and, inside a Cut \& Count formula, for N7; degree branching for A3) reduce the
parameter or the width before counting at the leaves.

\paragraph{Honest labels.} Table~\ref{tab:results} marks each result, after a literature search
of 28~October 2026 (papers, preprints, blogs and repositories). A1, A3, A3$'$ and N6 are known
consequences that we formalize. N3, N5 and N6$'$ are direct corollaries: the reduction is the
obvious one. A5a, A5a$'$ and A5b are superseded by concurrent work~\cite{Tul26}. The remaining
results, N2, N2$'$, N7, N7$'$, N7$''$, N8, N8$'$, N8$''$, A1$'$, the range of $k$ in A2 and the spectrum in A4, are ours: we did not find
them, which does not mean that they are certainly new.

\begin{table}[t]
\caption{The results. Time bounds hold for every $\varepsilon>0$ (resp.\ $\delta>0$); $C$ and
$n_0$ depend only on the stated parameters. $m$ is the number of edges, clauses or nonzero
entries.}\label{tab:results}
\centering\small
\begin{tabularx}{\textwidth}{@{}lXl@{}}
\toprule
& Problem and bound & Status\\
\midrule
N2 & $\oplus$connected vertex sets of size $\ell$ satisfying a CNF, any density; $C2^{\delta n}$ & ours\\
N2$'$ & Steiner tree, connected vertex cover, one-sided error, any density; $C2^{\varepsilon n}$ & ours\\
N7 & $\oplus$connected dominating sets of size $\ell$, any density; $C2^{\delta n}$ & ours\\
N7$'$ & connected dominating set, one-sided error, any density; $C2^{\varepsilon n}$ & ours\\
N7$''$ & Max Leaf Spanning Tree, minimum connected dominating set, one-sided error, any density; $C2^{\varepsilon n}$ & ours\\
N8 & $\oplus$connected vertex sets of size $\ell$ meeting $\le Dn$ given sets, any density; $C2^{\delta n}$ & ours\\
N8$'$ & Group Steiner tree, connected hitting set, connected total domination, one-sided error, any density; $C2^{\varepsilon n}$ & ours\\
N8$''$ & minimum Steiner tree, connected vertex cover, Group Steiner tree, one-sided error, any density; $C2^{\varepsilon n}$ & ours\\
A1 & \#CNF-SAT, any width, $m\le Dn$; $C2^{\delta n}$ & known~\cite{CIP06,anthropic}\\
A1$'$ & \#CNF-SAT, any width, $m\le n\,2^{(\lg n)^{1/4}/(\lg\lg n)^2}$, one machine; $2^{\delta n}$ & ours\\
A2 & \#$k$-colorings, all $k$ with $\lg k\le(\lg n)^{1/6-\theta}$, one machine; $2^{\varepsilon n}$ & ours (range of $k$)\\
N5 & \#homomorphisms $G\to H$, $H$ in the input, $\lg|V(H)|\le(\lg n)^{1/6-\theta}$; $2^{\varepsilon n}$ & corollary\\
A4 & Max-SAT spectrum, width $r$, $m\le Dn$; $C2^{\delta n}$ & ours (spectrum)\\
N3 & Potts density of states, Tutte polynomial on $(x-1)(y-1)=q$, $m\le Dn$; $C2^{\delta n}$ & corollary\\
N6 & 0/1 permanent; $C(n+2)^C2^{\varepsilon m}$ & known~\cite{openai}\\
N6$'$ & $\{-1,0,1\}$ permanent; $C(n+2)^C2^{\varepsilon m}$ & corollary\\
A5a & $\oplus$Hamiltonian cycles, degree $\le\Delta$; $2^{\varepsilon n}$ & superseded~\cite{Tul26}\\
A5a$'$ & $\oplus$Hamiltonian cycles, $m=o(n\sqrt{\lg n}/\lg\lg n)$, one machine; $2^{\varepsilon n}$ & superseded~\cite{Tul26}\\
A5b & Hamiltonian cycle, one-sided error, $m\le Dn$, one machine; $2^{\varepsilon n}$ & superseded~\cite{Tul26}\\
A3 & Vertex Cover; $C(n+2)^C2^{\varepsilon k}$ & known~\cite{IPZ01,CJ03}\\
A3$'$ & Nonblocker; $C(n+2)^C2^{\varepsilon k}$ & known~\cite{GdHKS16,DFFPR06}\\
\bottomrule
\end{tabularx}
\end{table}

\paragraph{Concurrent work.} Tulsyan~\cite{Tul26}, independently and also on top of Anthropic's
Lean development, proves that one machine computes the number of Hamiltonian cycles of a graph with
$n$ vertices and $m$ edges exactly and deterministically within $C2^{C(m+n)/\sqrt{\lg(n+2)}}$ steps,
and within $C2^{\delta n}$ steps for $m\le Dn$. His method uses Anthropic's network machinery
together with a matrix-tree identity, rather than the CNF counter. An exact count yields the parity
and a deterministic decision, so his result supersedes A5a, A5a$'$ and A5b. Our versions were
posted some hours earlier, on 27 October; we state them briefly in Section~\ref{sec:ham}. We rebuilt both theorems from his
repository on the same revision of Anthropic's development; they compile and use only the standard
axioms, and his count is the same quantity as our $\mathrm{hc}(G)$ of Section~\ref{sec:ham}. Since
then (29 October) he has added a Lean-verified exact algorithm for TSP with polynomial weights on
graphs with $m\le Dn$, a Lean-verified exact count in $2^{O(n\log\bar d/\sqrt{\lg n})}$ for average
degree $\bar d$, Lean proofs of decision and search for Hamiltonian cycles and of Hamiltonian paths,
and Lean-verified exact directed Hamiltonian cycle and ATSP. He also sketches deterministic
decision for Steiner tree, feedback vertex set, connected vertex cover, connected dominating set
and $k$-leaf spanning trees on graphs with $m\le Dn$ and bounded degree, explicitly as a conjecture that is checked only
numerically. He lists dominating set on dense graphs, and the exact counts of connected vertex
covers and connected dominating sets, as not covered. The two refutation papers list
corollaries of their own: OpenAI~\cite[Cor.~1.3]{openai} counts independent sets, $q$-colorings for
fixed $q$ and perfect matchings (this covers N6), and Anthropic~\cite[Cor.~7]{anthropic} counts
independent sets, cliques, vertex covers and dominating sets of each size, and $k$-colorings for
fixed $k$, on graphs of any density.

\paragraph{Related work.} The sparsification lemma~\cite{IPZ01} and the width--density duality of
Calabro, Impagliazzo and Paturi~\cite{CIP06} relate SAT variants with few variables; A1 is the latter
duality combined with the claimed algorithm. For the maximum in A4, the equivalence of Max-$r$-SAT
and Max-Cut with linearly many constraints with ETH is due to Jonsson, Lagerkvist, Schmidt and
Uppman~\cite[Thm.~8]{JLSU14}; Ganian, de Haan, Kanj and Szeider~\cite{GdHKS16} show the same for
Nonblocker and Dominating Set. For connectivity problems the direction we need was open:
Impagliazzo, Paturi and Zane list Hamiltonian circuit only as SNP-hard~\cite{IPZ01}, and Ganian et
al.\ ask whether Feedback Vertex Set reduces to 3-CNF-SAT in this way~\cite{GdHKS16}. Under ETH,
connected dominating set has no $2^{o(n)}$ algorithm even on intersection graphs of convex objects
in the plane~\cite{dBKB26}. For connected $G$ with $n\ge3$ vertices, $n$ minus the connected
domination number is the maximum number of leaves of a spanning tree~\cite{CWY00}. Dell, Husfeldt, Marx, Taslaman and
Wahl\'en~\cite{DHMTW14} prove \#ETH lower bounds for the permanent and the Tutte polynomial in
terms of the number of edges or nonzero entries; N3, N6 and N6$'$ show that, given the claimed
refutation, these problems become subexponential in the same parameter. Cut \&
Count~\cite{CNPPRW11} is the source of the parity reductions, and the isolation lemma of Calabro,
Impagliazzo, Kabanets and Paturi~\cite{CIKP08} is the source of the isolation step in
Section~\ref{sec:conn}. Cai and
Juedes~\cite{CJ03} show that $2^{o(k)}\poly(n)$ for Vertex Cover refutes ETH; A3 is the converse,
given the refutation. Cygan et al.~\cite{CFGKMPS17} show that, under ETH, deciding $G\to H$ needs
$|V(H)|^{\Omega(|V(G)|)}$ time in general. Our development uses Lean~4~\cite{lean4} and
Mathlib~\cite{mathlib}.

\paragraph{Organization.} Section~\ref{sec:prelim} describes the interface to Anthropic's
development. Section~\ref{sec:conn} treats connected solutions on graphs of any density,
Section~\ref{sec:counting} counting problems, Section~\ref{sec:ham} Hamiltonian cycles and
Section~\ref{sec:fpt} parameterized problems, each with algorithms, theorems and proof sketches.
Section~\ref{sec:formal} describes the formalization and the reference implementation, and
Section~\ref{sec:concl} discusses limitations.

%======================================================================
\section{Preliminaries and the interface}\label{sec:prelim}

A literal is a pair $(i,b)$ and holds under $x\in\{0,1\}^N$ iff $x_i=b$. A CNF is a list of clauses;
its width is the maximum clause length. We write $|x|$ for the number of true variables,
$\#\psi$ for the number of satisfying assignments of $\psi$ on $N$ variables, and $\cnt_s(\psi)$ for
the number of those with $|x|=s$. Logarithms $\lg$ are to base 2.

\paragraph{Machines and encodings.} A machine is a TM2 machine of Mathlib with finite stack
alphabets; its time is the number of steps until it halts with the stated output. Inputs are bit
strings. A graph on $n$ vertices is $1^n0$ followed by its adjacency matrix, a relation on $[h]$ is
its $h^2$ entries row by row, a number $t$ in unary is $1^t0$, and a CNF is $1^n0$ followed by its
clauses, each literal written in $\lfloor\lg 2n\rfloor+1$ bits.

\paragraph{The counters.} We use three theorems of Anthropic's development as black boxes.

\begin{theorem}[\SIZED{}, Anthropic]\label{thm:sized}
For every width $w$ and $\delta>0$ there are a machine $\SIZED_w$, given as a stepped computation
with polynomial-time steps, and a constant $c$, such that for all $N$, $s\le N$ and every CNF $\psi$
of width at most $w$ on $N$ variables, $\SIZED_w$ on $(1^N,s,\psi)$ returns $\cnt_s(\psi)$ within
$c\cdot2^{\delta N}$ steps of size at most $c(N+|\psi|+2)^c$.
\end{theorem}

\begin{theorem}[\COUNT{}, Anthropic's Main Theorem]\label{thm:count}
One machine \COUNT{} and a constant $C$: for every $w\ge1$ and every CNF $\psi$ of width $w$ with
$m\ge2$ clauses in which $N\ge4$ variables occur, given in a compact code of length $L$, \COUNT{}
outputs $\#\psi$ within $2^{C\min(wm/\sqrt{\lg m},\ w^2N\lg\lg N/\sqrt{\lg N})}+CL^C$ steps.
\end{theorem}

\begin{theorem}[\DOMCOUNT{}, Anthropic]\label{thm:dom}
For every $\delta>0$ there are a stepped machine \DOMCOUNT{} and $c$ that output the number of
dominating sets of size $s$ of a graph on $n$ vertices within $c\cdot2^{\delta n}$ steps.
\end{theorem}

The number of clauses enters Theorem~\ref{thm:sized} only through the polynomial size of a step.
Most of our reductions exploit this: their formulas have up to $n^2$ clauses, or as many as the input
has edges, but only linearly many variables.

\paragraph{Stepped computations and hosts.} Anthropic's machines are \emph{stepped}: a computation
is cut into polynomial-time steps, and a host theorem turns a stepped computation with $T$ steps of
state size $S$ into a TM2 machine with $q(T+1)(|x|+S+1)^q$ steps. A depth-first walk over the tree of
a step function that sums a stepped computation over the leaves is again stepped
(\lean{walkSum}); its number of steps is bounded via the number of leaves. We add one host.

\begin{lemma}[FPT host, \lean{tm_fpt}]\label{lem:host}
Let a stepped computation run on inputs $x$ of size $n(x)$ with parameter $p(x)$ in at most
$c\,2^{\delta p(x)}$ steps of state size $a(n+2)^e$. Then there are $C$ and one TM2 machine with
finite alphabets that produces the same output within $\lfloor C(n+2)^C2^{\delta p(x)}\rfloor$ steps.
\end{lemma}

\paragraph{The pattern.} Every algorithm below has the form
\[
x\ \xrightarrow{\ \text{front end, poly time}\ }\ (N,s,\psi)\ \xrightarrow{\ \SIZED\text{ or }\COUNT\ }\ \text{count}\ \xrightarrow{\ \text{post-step}\ }\ \text{answer},
\]
or a tree walk over such leaves. The front end is a polynomial-time function, proved so in Lean
(\lean{PolyT}), and it is joined to the counter by \lean{Stepped.mapIn} and \lean{Stepped.mapOut}.
The proof obligations are always the same: the formula has the claimed width and number of
variables, and its (sized) count is the quantity in the statement.

%======================================================================
\section{Connected solutions on graphs of any density}\label{sec:conn}

Cut \& Count~\cite{CNPPRW11} replaces connectivity by the parity of a count. A vertex set $X$ has
$2^{\mathrm{cc}(G[X])}$ maps $s:X\to\{0,1\}$ that are constant on the components of $G[X]$, and
$2^{\mathrm{cc}(G[X])-1}$ is odd iff $G[X]$ is connected. The side bits add $n$ variables, and the
edges of $G$ enter only as clauses. Theorem~\ref{thm:sized} is insensitive to the number of clauses,
so the graph may be dense.

\subsection{The parity of connected solutions (N2)}

\begin{algorithm}[t]
\caption{Parity of connected solutions (N2). Lean: \lean{Conn.connF}, \lean{ConnTM.connM}.}\label{alg:conn}
\begin{algorithmic}[1]
\Require graph $G$ on $n$ vertices, $1\le\ell\le n$, CNF $\varphi$ of width $r$ on the vertex bits
\State $x_v$: $v\in X$; \ $s_v=x_{n+v}$: the side of $v$; \ $\bar s_v=x_{2n+v}$
\State $\psi\gets\varphi$, $\ \bar s_v=\neg s_v$, $\ x_v\vee\neg s_v$, and for each edge $uv$: $x_u\wedge x_v\to(s_u=s_v)$
\State $c\gets\SIZED_{r+4}(3n,\ n+\ell,\ \psi)$ \Comment{width $\le\max(r,4)$; $c=\sum_X2^{\mathrm{cc}(G[X])}$}
\State \Return $\lfloor c/2\rfloor\bmod2$
\end{algorithmic}
\end{algorithm}

\begin{theorem}[\lean{conn_parity_tm}, \lean{cvc_parity_tm}]\label{thm:conn}
For every width $r$, exponent $D$ and $\delta>0$, one machine, given a graph $G$ on $n$ vertices of
any density, $1\le\ell\le n$ and a CNF $\varphi$ on the vertices with clauses of at most $r$ literals
and at most $(n+2)^D$ clauses, outputs the parity of the number of $X\subseteq V$ with $|X|=\ell$,
$X\models\varphi$ and $G[X]$ connected, within $C2^{\delta n}$ steps.
\end{theorem}

\begin{proof}[Proof sketch]
Each pair $(s_v,\bar s_v)$ has exactly one true bit, so the solutions of size $n+\ell$ are the pairs
$(X,s)$ with $|X|=\ell$, $X\models\varphi$ and $s$ constant on the components of $G[X]$ and $0$
outside $X$ (\lean{cnt_stepA}, \lean{cnt_stepB}). There are $2^{\mathrm{cc}(G[X])}$ of them for each
$X$, and $2^{\mathrm{cc}-1}$ is odd iff $G[X]$ is connected (\lean{sig_parity}). The formula has $3n$
variables whatever the number of edges.
\end{proof}

With $\varphi$ the unit clauses $x_t$ for a terminal set $T$, the sets are the vertex sets of Steiner
trees on $\ell$ vertices; with $\varphi=\{x_u\vee x_v: uv\in E\}$, the connected vertex covers of size
$\ell$ (a separate machine on input $G,\ell$).

\subsection{Isolation for arbitrary sets}

\begin{theorem}[\lean{isolation}; after~\cite{CIKP08}]\label{thm:iso}
For every $\delta>0$ there are $k,t$ such that for all $M\ge M_0(\delta)$ and every nonempty
$S\subseteq\{0,1\}^M$, at least a $2^{-\delta M}$ fraction of the seeds leaves exactly one point of
$S$. A seed consists of $3M+3$ parity checks on $k$ coordinates each (with right-hand sides), a
number of them to use, and for every coordinate a value in $\{0,\dots,t-1\}$ and a bit; the
coordinate is fixed to the bit when the value is 0.
\end{theorem}

\begin{proof}[Proof sketch]
One check on $k$ coordinates fails to separate two points at distance $d$ with probability
$(1+(1-d/M)^k)/2$ (\lean{two_mul_card_rv}, by flipping one coefficient). Call $x\in S$ semi-isolated if no point of $S$ at
distance $\ge\rho$ passes the same first $m$ checks. For the least $m$ with $2|S|q^m\le1$, half of the
check arrays make a given $x$ semi-isolated (\lean{card_not_semi}); the semi-isolated points with the
same check values lie in one Hamming ball (\lean{card_fiber_le}). If the survivors lie within
distance $<2\rho$ of each other, fixing the coordinates where the two furthest differ leaves one point
(\lean{stepII}). The count \lean{iso_count} is exact; \lean{iso_real} takes $\rho=\lfloor
M/a\rfloor+1$ and $k=ab$, and the theorem picks $t$, $b$, $a$ against $\delta$.
\end{proof}

\cite{CIKP08} state the lemma for solution sets of $k$-CNFs; its proof uses only $S\neq\emptyset$,
and that is what we formalize. A check on $k$ coordinates is a chain of $k$ auxiliary variables
$a_l=a_{l-1}\oplus(c_l\wedge x_{j_l})$ in clauses of width $\le3$, and the chain bits are determined by
the main variables. So isolating among the solutions of a CNF adds $O(Mk)$ variables, linear in $M$.

\subsection{\texorpdfstring{Steiner tree and connected vertex cover with one-sided error (N2$'$)}{Steiner tree and connected vertex cover with one-sided error (N2')}}

\begin{algorithm}[t]
\caption{Connected solutions with one-sided error (N2$'$). Lean: \lean{ConnI.ciF}, \lean{ConnITM.preR}.}\label{alg:connrand}
\begin{algorithmic}[1]
\Require $G$, $\ell$, $\varphi$ as in Algorithm~\ref{alg:conn}; a seed $\omega$ with $3n+3$ checks on the $n$ vertex bits
\State $\psi\gets$ the formula of Algorithm~\ref{alg:conn} on $3n$ variables; \ $z=x_{3n}$, add $\neg z$
\For{each check $i<3n+3$ with terms $(j_0,c_0),\dots,(j_{k-1},c_{k-1})$}
  \State add $a_{i,l}=a_{i,l-1}\oplus(c_l\wedge x_{j_l})$ for $l<k$, where $a_{i,l}=x_{3n+1+ik+l}$ and $a_{i,-1}=z$
\EndFor
\State add $a_{i,k-1}=$ (right-hand side of check $i$) for the checks in use, and $x_j=$ bit$_j$ for the fixed coordinates
\State for every chain bit $a$, a complement bit $\bar a=\neg a$ \Comment{keeps the target size fixed}
\State $c\gets\SIZED_{r+4}(3n+1+2(3n+3)k,\ n+\ell+(3n+3)k,\ \psi)$
\State \Return $\lfloor c/2\rfloor\bmod2$
\end{algorithmic}
\end{algorithm}

\begin{theorem}[\lean{conn_randomized}]\label{thm:connrand}
For every $r$, $D$, $\delta>0$ and $\varepsilon>0$ there are $k,t,n_0,C$ and one machine with the
following property. For every graph $G$ on $n\ge n_0$ vertices of any density, $1\le\ell\le n$,
every CNF $\varphi$ as in Theorem~\ref{thm:conn} and every seed $\omega$, the machine halts on
$(G,\ell,\omega,\varphi)$ within $C2^{\varepsilon n}$ steps with a bit $\mathrm{out}(\omega)$. If no
vertex set of size $\ell$ satisfies $\varphi$ and induces a connected subgraph,
$\mathrm{out}(\omega)=0$ for every seed; if one does, $\mathrm{out}(\omega)=1$ for at least a
$2^{-\delta n}$ fraction of the seeds. In particular this decides with one-sided error whether $G$
has a Steiner tree on $\ell$ vertices and whether it has a connected vertex cover of size $\ell$.
\end{theorem}

\begin{proof}[Proof sketch]
Algorithm~\ref{alg:connrand}. A chain bit and its complement hold exactly one true bit, so the
solutions with $n+\ell+(3n+3)k$ true variables are those of Algorithm~\ref{alg:conn} with $n+\ell$
true variables whose vertex set passes the checks and the fixings (\lean{stepC}--\lean{stepE}), and half their
number has the parity of the number of connected solutions of size $\ell$ that pass them
(\lean{ci_parity}). If there is no connected solution, this parity is 0 for every seed. If there is
one, Theorem~\ref{thm:iso} applied to the set of connected solutions in $\{0,1\}^n$ leaves exactly one
for a $2^{-\delta n}$ fraction of the seeds, and there the parity is 1 (\lean{ci_decide}). The
formula has $3n+1+2(3n+3)k=O(n)$ variables, whatever the number of edges.
\end{proof}

Running the machine on $R$ independent seeds and answering yes if some run outputs 1 never errs on
a no-instance. On a yes-instance it errs on at most a fraction $e^{-R2^{-\delta n}}$ of the
$R$-tuples of seeds, which is at most $e^{-2^{\delta n}}$ for $R\ge2^{2\delta n}$
(\lean{conn_amplified}).

\subsection{\texorpdfstring{Connected dominating set and Max Leaf (N7, N7$'$, N7$''$)}{Connected dominating set and Max Leaf (N7, N7', N7'')}}

A connected dominating set is a set $X$ with $G[X]$ connected such that every vertex is in $X$ or
has a neighbour in $X$. Domination is the CNF $\varphi_{\mathrm{dom}}=\{\bigvee_{u\in N[v]}x_u : v\in
V\}$, but its clauses have unbounded width, so Theorem~\ref{thm:conn} does not apply. Anthropic's
algorithm for dominating sets~\cite[Cor.~7]{anthropic} (Theorem~\ref{thm:dom}) shortens these
clauses with a search tree after Calabro, Impagliazzo and Paturi~\cite{CIP06}: in a clause with more
than $w$ variables, let $H$ be its first $w$ variables; in the \emph{light} child the clause is
replaced by $\bigvee H$, in the \emph{heavy} child the variables of $H$ are false and are deleted
from all open clauses. We run this tree with the clauses of domination open and the clauses of
Cut \& Count closed.

\begin{algorithm}[t]
\caption{Parity of connected dominating sets (N7). Lean: \lean{CdsTree.rootN}, \lean{CdsTM.cdsM}.}\label{alg:cds}
\begin{algorithmic}[1]
\Require graph $G$ on $n$ vertices, $1\le\ell\le n$
\State open $\gets\{N[v]: v\in V\}$; \ closed $\gets$ the clauses of Algorithm~\ref{alg:conn} for $\varphi=\emptyset$, on $3n$ variables
\State $c\gets0$
\For{each leaf (open$'$, closed$'$) of the tree of the shortening step with width $w$}
  \State $c\gets c+\SIZED_{w+4}(3n,\ n+\ell,\ \text{open}'\text{ as positive clauses}\cup\text{closed}')$
\EndFor
\State \Return $\lfloor c/2\rfloor\bmod2$
\end{algorithmic}
\end{algorithm}

\begin{theorem}[\lean{cds_parity_tm}]\label{thm:cds}
For every $\delta>0$, one machine, given a graph $G$ on $n$ vertices of any density and
$1\le\ell\le n$, outputs the parity of the number of connected dominating sets of $G$ with $\ell$
vertices, within $C2^{\delta n}$ steps.
\end{theorem}

\begin{proof}[Proof sketch]
Every assignment satisfies exactly one child of a node, so the sized counts of the leaves add up to
that of the root (\lean{cntS_step}, \lean{sum_leaves}), which is the formula of
Algorithm~\ref{alg:conn} with $\varphi=\varphi_{\mathrm{dom}}$. A light step removes an open clause
and a heavy step removes $w$ of the variables that occur in open clauses. Anthropic's invariant
bounds these variables by the number of variables of the formula; ours (\lean{InvC}) bounds them
by $n$ while the formula has $3n$ variables and arbitrary closed clauses of width at most $w+4$. So
a path has at most $n$ light and $n/w$ heavy steps, the tree has at most
$\sum_{j\le2n/w}\binom{2n}{j}\le2^{\varepsilon n}$ leaves for $w$ large enough, and every leaf has
width at most $w+4$ (\lean{shortensN}). With $\varepsilon=\delta/4$ and $\SIZED_{w+4}$ at $\delta/12$
on $3n$ variables, the walk takes $C2^{\delta n/2}$ steps; Theorem~\ref{thm:conn}'s parity argument
(\lean{conn_parity}) gives the output.
\end{proof}

\begin{theorem}[\lean{cds_randomized}]\label{thm:cdsrand}
For every $\delta>0$ and $\varepsilon>0$ there are $k,t,n_0,C$ and one machine that, for every
graph $G$ on $n\ge n_0$ vertices of any density, $1\le\ell\le n$ and every seed $\omega$ as in
Theorem~\ref{thm:connrand}, halts on $(G,\ell,\omega)$ within $C2^{\varepsilon n}$ steps with a bit
$\mathrm{out}(\omega)$. If $G$ has no connected dominating set of size $\ell$, $\mathrm{out}(\omega)=0$
for every seed; if it has one, $\mathrm{out}(\omega)=1$ for at least a $2^{-\delta n}$ fraction of the
seeds.
\end{theorem}

\begin{proof}[Proof sketch]
The closed clauses at the root are now those of Algorithm~\ref{alg:connrand} for $\varphi=\emptyset$,
on $3n+1+2(3n+3)k$ variables; the open ones are still the $n$ clauses of domination on the vertex
bits, so the tree and its bound are those of Theorem~\ref{thm:cds}. The sum over the leaves is the
count of Algorithm~\ref{alg:connrand} for $\varphi=\varphi_{\mathrm{dom}}$, and \lean{ci_decide} gives
the one-sided error. Amplification is as for N2$'$ (\lean{cds_amplified}).
\end{proof}

\begin{theorem}[\lean{maxleaf_randomized}, \lean{mincds_randomized}]\label{thm:maxleaf}
For every $\delta>0$ and $\varepsilon>0$ there are $k,t,n_0,C$ and one machine that, for every
graph $G$ on $n\ge n_0$ vertices of any density, every $1\le L<n$ and every seed $\omega$, halts on
$(G,n-L,\omega)$ within $C2^{\varepsilon n}$ steps with a bit $\mathrm{out}(\omega)$. If $G$ has no
spanning tree with at least $L$ leaves, $\mathrm{out}(\omega)=0$ for every seed; if it has one,
$\mathrm{out}(\omega)=1$ for at least a $2^{-\delta n}$ fraction of the seeds. On $(G,\ell,\omega)$
with $1\le\ell\le n$ the same machine decides, in the same sense, whether $G$ has a connected
dominating set with at most $\ell$ vertices.
\end{theorem}

\begin{proof}[Proof sketch]
It is the machine of Theorem~\ref{thm:cdsrand}. Two facts about graphs, proved with Mathlib's walks
and trees, translate the conditions. First, if $X$ is a connected dominating set with
$|X|<\ell\le n$, some vertex outside $X$ has a neighbour in $X$, and adding it keeps $X$ connected
and dominating; so a connected dominating set with at most $\ell$ vertices exists iff one with
exactly $\ell$ does (\lean{cds_le_iff}). Second, for $n\ge3$ and $1\le L<n$, $G$ has a spanning
tree with at least $L$ leaves iff it has a connected dominating set with $n-L$
vertices~\cite{CWY00} (\lean{leaf_iff}). In a spanning tree on $n\ge3$ vertices the non-leaves
induce a connected subgraph, because the inner vertices of a path have degree at least 2, and they
dominate, because the neighbour of a leaf is not a leaf. Conversely, a spanning tree of $G[X]$
together with one edge from each vertex outside $X$ to a neighbour in $X$ is a spanning tree of
$G$ in which every vertex outside $X$ is a leaf. A disconnected graph has neither, so no test of
connectivity is needed.
\end{proof}

The proof of Theorem~\ref{thm:cds} uses only that the open clauses are positive, have no repeated
variables and live on the vertex bits. With the invariant bounding the number of open clauses by
$Dn$ and their variables by $n$ separately, as for A1, a path of the tree has at most $Dn$ light and
$n/w$ heavy steps, and the tree of depth $(D+1)n$ has at most
$\sum_{j\le2n/w}\binom{(D+2)n}{j}\le2^{\varepsilon n}$ leaves for a suitable $w$
(\lean{shortensP}).

\begin{theorem}[\lean{pos_parity_tm}, \lean{pos_randomized}]\label{thm:pos}
For every $D$ and $\delta>0$ there are $C$ and one machine that, given a graph $G$ on $n$ vertices of
any density, $1\le\ell\le n$ and a CNF $\varphi$ on the vertices with only positive literals, at most
$n$ literals per clause and at most $Dn$ clauses of any width, outputs within $C2^{\delta n}$ steps
the parity of the number of sets $X$ with $|X|=\ell$ that meet every clause of $\varphi$ and induce a
connected subgraph. For every $D$, $\delta>0$ and $\varepsilon>0$ there are $k,t,n_0,C$ and one
machine that, for $n\ge n_0$ and every seed $\omega$ as in Theorem~\ref{thm:connrand}, halts on
$(G,\ell,\omega,\varphi)$ within $C2^{\varepsilon n}$ steps with a bit $\mathrm{out}(\omega)$, which is
$0$ for every seed if no such $X$ exists and $1$ for at least a $2^{-\delta n}$ fraction of the seeds if
one does.
\end{theorem}

\begin{proof}[Proof sketch]
The front end keeps the variables of each clause without repeats; a positive clause has the same
models as the clause of its variables (\lean{cntS_opens}). The rest is the proof of
Theorems~\ref{thm:cds} and~\ref{thm:cdsrand} with $\varphi$ in place of $\varphi_{\mathrm{dom}}$.
\end{proof}

With $\varphi$ the groups, the sets $X$ are the vertex sets of the group Steiner trees with $\ell$
vertices; with $\varphi$ a set system, the connected hitting sets of size $\ell$; with $\varphi$ the open
neighbourhoods, the connected total dominating sets of size $\ell$. In each case at most $Dn$ sets are
allowed. Theorem~\ref{thm:conn} allows negative literals, but only in clauses of bounded width.

For a positive $\varphi$, adding a neighbour to a connected set that meets $\varphi$ gives another one,
and in a connected graph some vertex outside a set $X\ne\emptyset,V$ has a neighbour in $X$. So a
solution with at most $\ell\le n$ vertices grows to one with exactly $\ell$ (\lean{conn_grow},
\lean{min_iff}), and on connected graphs the machines of Theorems~\ref{thm:connrand}
and~\ref{thm:pos}, run with $\ell$, decide with one-sided error whether a solution with at most
$\ell$ vertices exists (N8$''$, \lean{min_conn_randomized}, \lean{min_pos_randomized}): minimum
Steiner tree (in vertices), minimum connected vertex cover and minimum Group Steiner tree.

%======================================================================
\section{Counting}\label{sec:counting}

\subsection{\texorpdfstring{CNF of any width (A1, A1$'$)}{CNF of any width (A1, A1')}}

\begin{algorithm}[t]
\caption{\#CNF-SAT of any width (A1). Lean: \lean{CnfTM.cnfM}.}\label{alg:cnf}
\begin{algorithmic}[1]
\Require CNF $\varphi$ on $n$ variables, width bound $w$
\State root: variables $x_{2u}$ (``$u$ true''), $x_{2u+1}$ (``$u$ false''); closed clauses
$x_{2u}\vee x_{2u+1}$, $\neg x_{2u}\vee\neg x_{2u+1}$; open clauses: each clause of $\varphi$ as a set
of positive variables
\State $\mathit{total}\gets0$
\For{each node of the tree, depth first}
  \If{an open clause $K$ has more than $w$ variables; $H$ := its first $w$}
    \State light child: remove $K$ from the open clauses, add $\bigvee_{i\in H}x_i$ to the closed ones
    \State heavy child: add $\neg x_i$ for $i\in H$, delete $H$ from every open clause
  \Else\ $\mathit{total}\gets\mathit{total}+\SIZED_w(2n,\,n,\,\text{open}\cup\text{closed})$
  \EndIf
\EndFor
\State \Return $\mathit{total}$
\end{algorithmic}
\end{algorithm}

\begin{theorem}[\lean{cnf_counts}]
For every $D$ and $\delta>0$, one machine counts the satisfying assignments of every CNF with $n$
variables and at most $Dn$ clauses, of any width and with no literal repeated within a clause, within
$C2^{\delta n}$ steps.
\end{theorem}

\begin{proof}[Proof sketch]
Algorithm~\ref{alg:cnf}. Since $|x|=n$ forces exactly one of $x_{2u},x_{2u+1}$, the solutions of size
$n$ at the root are the assignments of $\varphi$. A node's count is the sum of its children's
(either some variable of $H$ is true, or all are false), and every leaf has width $\le w$. The tree is
Anthropic's clause-shortening tree from their Dominating Set corollary; we replace its invariant,
which needs at most $n$ clauses, by one for $Dn$ clauses. The tree has at most
$2^{(2/w)\lg(ew(D+2))n}$ leaves, which is $2^{\delta n/2}$ for $w$ large.
\end{proof}

\begin{theorem}[\lean{cnf_grow}]
One machine: for every $\delta>0$ there is $n_0$ such that on every CNF with $n\ge n_0$ variables and
at most $n\cdot 2^{(\lg n)^{1/4}/(\lg\lg n)^2}$ clauses of any width, with no literal repeated
within a clause, it outputs the number of satisfying assignments within $2^{\delta n}$ steps.
\end{theorem}

\begin{proof}[Proof sketch]
Algorithm~\ref{alg:cnf} with the width read off the node: for $N$ variables, $\ell=\lfloor\lg
N\rfloor$ and $L=\lfloor\lg\ell\rfloor$, take $w=2^{\lfloor L/4\rfloor-\lfloor\lg L\rfloor}+1\approx
\ell^{1/4}/L$, and count each leaf with \COUNT. Its exponent $Cw^2N\lg\lg N/\sqrt{\lg N}$ is at most
$32CN/L=o(n)$, since $w^2L^2\le16\sqrt\ell$. For $D\le2^s+1$ with $s=(\lg n)^{1/4}/(\lg\lg n)^2$, so
that $s\le4w/(L-1)$, the coefficient $(2/w)\lg(ew(D+2))$ of the tree bound is
at most $(8+6\sqrt w)/w+8/(L-1)\to0$.
\end{proof}

\subsection{Colorings and homomorphisms (A2, N5)}

Let $H$ be a relation on $[h]=\{0,\dots,h-1\}$ (loops allowed, not necessarily symmetric) and
$\hom(G,H)$ the number of maps $f:V(G)\to[h]$ with $H(f(u),f(v))$ for every edge $uv$. For
$H=K_k$ these are the proper $k$-colorings.

\begin{algorithm}[t]
\caption{Homomorphisms $G\to H$ (N5; A2 is $H=K_k$). Lean: \lean{Hom.homB}, \lean{colB}.}\label{alg:hom}
\begin{algorithmic}[1]
\Require graph $G$ on $n$ vertices, relation $H$ on $[h]$
\State $b\gets\lceil\lg h\rceil$; the bits $x_{vb},\dots,x_{vb+b-1}$ spell the image of $v$
\State for each $v$ and each code $h\le c<2^b$: clause ``the bits of $v$ do not spell $c$''
\State for each edge $uv$ (both orders) and each $(a,c)$ with $\neg H(a,c)$: clause ``not ($u$ spells $a$ and $v$ spells $c$)''
\State tautologies $x_i\vee\neg x_i$ for $i<nb$; unit clauses $x_{nb},\dots,x_{nb+3}$
\State \Return $\COUNT(nb+4,\psi)$
\end{algorithmic}
\end{algorithm}

\begin{theorem}[\lean{colorings_subexp}, \lean{hom_subexp}, \lean{hom_subexp_graph}]
One machine, given $h$ in unary, the matrix of $H$ and an $n$-vertex graph $G$, outputs $\hom(G,H)$
within $2^{\varepsilon n}$ steps, for every $\varepsilon>0$, $0<\theta<1/6$, $n\ge
n_0(\theta,\varepsilon)$ and all $h\ge2$ with $\lg h\le(\lg n)^{1/6-\theta}$. For a simple graph $H$
the output is the number of Mathlib homomorphisms $G\to_g H$; for $H=K_k$ it is the number of proper
$k$-colorings.
\end{theorem}

\begin{proof}[Proof sketch]
The satisfying assignments of the formula of Algorithm~\ref{alg:hom} are exactly the homomorphisms
(\lean{cntN_homB}). It has $N=nb+4$ variables and width $2b$, so the exponent of
Theorem~\ref{thm:count} is $C\cdot4b^2\cdot N\lg\lg N/\sqrt{\lg N}=O(n(\lg n)^{-2\theta})$ for
$b\le2(\lg n)^{1/6-\theta}$, because $b^3\le8(\lg n)^{1/2-3\theta}$ and $\lg\lg N\le(\lg
N)^\theta/(\theta\ln2)$. The at most $n^2h^2+O(nh)$ clauses only enter the polynomial factor.
\end{proof}

Taking $h=2$ and $H$ an edge with a loop at one end counts the independent sets of $G$ in Mathlib's sense, for graphs of any density
(\lean{numHom_hInd}, \lean{indep_subexp}).

\subsection{The Max-SAT spectrum (A4) and the Potts model (N3)}

\begin{theorem}[\lean{maxsat_spectrum}]
For every width $r$, density $D$ and $\delta>0$, one machine outputs, given a CNF $\varphi$ with
$m\le Dn$ clauses of at most $r$ literals and $t\le m$, the number of assignments that satisfy
exactly $t$ clauses, within $C2^{\delta n}$ steps.
\end{theorem}

\begin{proof}[Proof sketch]
Use the $2n$ variables of A1 and one variable $y_j=x_{2n+j}$ per clause $K_j$, with clauses
$y_j\leftrightarrow K_j$ (width $r+1$), the literal $(u,b)$ read on $x_{2u}$. A solution with $n+t$
true variables is an assignment of $\varphi$ together with the indicator of its $t$ satisfied
clauses (\lean{cntS_specF}). Return $\SIZED_{r+2}(2n+m,\,n+t,\,\psi)$.
\end{proof}

For graphs with $O(n)$ edges, the Max-Cut spectrum and the density of states of the $\pm1$ Ising
model follow from N3 below with $q=2$ (\lean{maxcut_spectrum}, \lean{ising_spectrum}). For a graph $G$ and $q$ colors let $N_j$ be the number of colorings $\sigma:V\to[q]$
with exactly $j$ monochromatic edges, $k(A)$ the number of components of $(V,A)$, and
$T(G;x,y)=\sum_{A\subseteq E}(x-1)^{k(A)-k(E)}(y-1)^{|A|+k(A)-n}$.

\begin{theorem}[\lean{potts_spectrum}, \lean{potts_fk}, \lean{tutte_hyperbola}, \lean{tutte_tm}]
For every $q\ge1$, $D$ and $\delta>0$, one machine outputs $N_j$, given $G$ with $m\le Dn$ edges
and $j\le m$, within $C2^{\delta n}$ steps. In every commutative ring
$\sum_jN_j(1+v)^j=\sum_{A\subseteq E}q^{k(A)}v^{|A|}$, and in every field, if $(x-1)(y-1)=q$ and
$x,y\ne1$, then $(x-1)^{k(E)}(y-1)^nT(G;x,y)=\sum_jN_jy^j$.
\end{theorem}

\begin{proof}[Proof sketch]
Variables $x_{vq+c}$ (``$v$ has color $c$'') and $y_i=x_{nq+i}$ (``edge $i$ is monochromatic''),
$nq+m$ in all. Clauses: each vertex has at least one color (width $q$) and at most one (width 2); for
each edge $(a,b)$ and color $c$, $x_{aq+c}\wedge x_{bq+c}\to y_i$ and $y_i\wedge x_{aq+c}\to
x_{bq+c}$. The solutions with $n+j$ true variables are the colorings with $j$ monochromatic edges
(\lean{cnt_pottsF}); return $\SIZED_{q+3}(nq+m,\,n+j,\,\psi)$. The identities are the
Fortuin--Kasteleyn expansion~\cite{FK72}: expand $\prod_e(1+v[e\text{ mono}])$ over $A\subseteq E$;
the colorings constant on the components of $(V,A)$ are $q^{k(A)}$ many (\lean{card_const}).
\end{proof}

For $q=2$, the colorings with $j$ monochromatic edges are the vertex sets whose cut has $m-j$ edges
(\lean{pottsSpec_two}), and the spin configurations with $\sum_{uv\in E}s_us_v=2j-m$
(\lean{ising_spectrum}). $N_0$ is the number of proper $q$-colorings in Mathlib's sense
(\lean{pottsSpec_zero}). It is the chromatic polynomial at $q$, and
$N_0=(-q)^{k(E)}(-1)^nT(G;1-q,0)$ (\lean{chromatic_tutte}). At the point $(0,1-q)$ the theorem gives $T(G;0,1-q)$ (\lean{flow_tutte}),
which by Tutte's theorem~\cite{Tutte54} is $(-1)^{m-n+k(E)}$ times the number of nowhere-zero
$\mathbb Z_q$-flows. Under \#ETH, \cite{DHMTW14} prove lower bounds of the form
$2^{\Omega(n/\mathrm{polylog}\,n)}$ for such evaluations on sparse graphs.

\subsection{\texorpdfstring{The permanent (N6, N6$'$)}{The permanent (N6, N6')}}

\begin{algorithm}[t]
\caption{The $\{-1,0,1\}$ permanent (N6$'$; N6 returns $\SIZED_2(m,n,\psi)$ after line~2). Lean:
\lean{Perm.permF}, \lean{SPerm.chainF}, \lean{SPermTM.spermM}.}\label{alg:perm}
\begin{algorithmic}[1]
\Require $A,S\in\{0,1\}^{n\times n}$: $M_{ij}=0$ if $\neg A_{ij}$, $-1$ if $A_{ij}\wedge S_{ij}$, $1$ otherwise
\State $e_0,\dots,e_{m-1}\gets$ the entries with $A_{ij}=1$, row by row; one variable $x_t$ each;
\textbf{if} $m<n$ \textbf{return} 0
\State $\psi\gets\{\neg x_s\vee\neg x_t : s<t,\ e_s,e_t\text{ share a row or a column}\}$
\Comment{width 2}
\State $g_0,\dots,g_{r-1}\gets$ the indices $t$ with $S$ at $e_t$
\For{$b\in\{0,1\}$}
  \State $\psi_b\gets\psi\cup\{\neg z\}\cup\{p_k=p_{k-1}\oplus x_{g_k},\ q_k=\neg p_k : k<r\}\cup\{p_{r-1}=b\}$
  \Statex \hspace{\algorithmicindent} where $z=x_m$, $p_{-1}=z$, $p_k=x_{m+1+k}$, $q_k=x_{m+1+r+k}$
  \Comment{width 3}
  \State $N_b\gets\SIZED_3(m+1+2r,\ n+r,\ \psi_b)$
\EndFor
\State \Return $N_0-N_1$
\end{algorithmic}
\end{algorithm}

\begin{theorem}[\lean{perm_fpt}, \lean{sperm_fpt}, \lean{permanent_signed}]\label{thm:perm}
For every $\varepsilon>0$ there are $C$ and one machine that, given a bit $b$, $n$ in unary and
$A,S\in\{0,1\}^{n\times n}$, outputs
$N_b=\#\{\sigma\in S_n : \forall i\ A_{\sigma(i),i}=1,\ |\{i: S_{\sigma(i),i}=1\}|\equiv b\pmod 2\}$
within $C(n+2)^C2^{\varepsilon m}$ steps, where $m$ is the number of ones of $A$. The permanent of
$M$ (Mathlib's \lean{Matrix.permanent} over $\mathbb Z$) equals $N_0-N_1$. For $S=0$ a separate
machine outputs the permanent of the 0/1 matrix $A$ directly, in the same time.
\end{theorem}

\begin{proof}[Proof sketch]
An assignment of $\psi$ with exactly $n$ true variables chooses $n$ ones of $A$ with no two in a row
or a column, i.e.\ a permutation $\sigma$ with $A_{\sigma(i),i}=1$ for all $i$; the map is a
bijection (\lean{cntS_permF}), and these $\sigma$ are the nonzero terms of the permanent of the 0/1
matrix (\lean{permanent_eq}). In $\psi_b$ the chain bits are determined by the choice
(\lean{chain_det}), $p_k$ is the parity of the first $k+1$ entries $-1$ that are chosen
(\lean{chainVal_eq}), and every pair $(p_k,q_k)$ has exactly one true bit. So the solutions of size
$n+r$ are the permutations of parity $b$ (\lean{cnt_sperm}). The term of $\sigma$ in the permanent
of $M$ is $(-1)^{j(\sigma)}$, where $j(\sigma)$ is the number of entries $-1$ on $\sigma$ (\lean{permanent_signed}). The formula has $m+1+2r\le3m+1$
variables, so $\SIZED_3$ at $\delta=\varepsilon/3$ takes $c\,2^{\varepsilon/3}2^{\varepsilon m}$
steps, and Lemma~\ref{lem:host} gives the machine.
\end{proof}

Under \#ETH the permanent of $\{-1,0,1\}$ matrices has no $\exp(o(m))$ algorithm, and that of 0/1
matrices none in $\exp(o(m/\log m))$~\cite[Thm.~1.3]{DHMTW14}. Ryser's formula~\cite{Ryser63} takes
$O(2^nn)$ time.

%======================================================================
\section{Hamiltonian cycles}\label{sec:ham}

Let $\mathrm{hc}(G)$ be the number of connected spanning subgraphs of $G$ with all degrees 2. In
Section~\ref{sec:formal} we prove that it is the number of Hamiltonian cycles in Mathlib's sense.
Tulsyan~\cite{Tul26} computes $\mathrm{hc}(G)$ exactly and deterministically, which supersedes the
results of this section; we obtained them independently and state them briefly. The encoding is
that of Section~\ref{sec:conn} with edge bits in place of vertex bits.

\begin{algorithm}[t]
\caption{Parity of the number of Hamiltonian cycles (A5a, A5a$'$). Lean: \lean{Ham.hcF}, \lean{HamS.hcS}.}\label{alg:hc}
\begin{algorithmic}[1]
\Require graph $G$ on $n$ vertices with edges $e_0,\dots,e_{m-1}$
\State $x_i$: edge $e_i$ is chosen; $s_v=x_{m+v}$: the side of $v$
\State for each $v$: exactly two chosen edges at $v$
\Statex \hspace{\algorithmicindent}A5a: $\neg x_a\vee\neg x_b\vee\neg x_c$ for distinct edges at $v$, and ``at least two'' (width $\deg v$)
\Statex \hspace{\algorithmicindent}A5a$'$: a sequential counter along the edges at $v$, two bits per (edge, end), width 3
\State for each edge $e_i=ab$: $x_i\to(s_a=s_b)$; \ $\neg s_0$; \ $x_i\vee\neg x_i$ for every variable
\State \Return $\COUNT(N,\psi)\bmod2$ \Comment{$N=m+n$ (A5a), $N=5m+n+2$ (A5a$'$)}
\end{algorithmic}
\end{algorithm}

\begin{theorem}[\lean{hc_parity_subexp}, \lean{hc_parity_little_o}]\leavevmode\\
For each $\Delta$, one machine decides whether a graph with all degrees at most $\Delta$ has an odd
number of Hamiltonian cycles within $2^{\varepsilon n}$ steps for $n\ge n_0(\varepsilon)$. One machine
decides it for all graphs of any degree with at most $n\,g(n)$ edges, within $2^{\varepsilon n}$ steps
for $n\ge n_0(g,\varepsilon)$, for every $\varepsilon>0$ and every $g$ with $g(n)\,r(n)\to0$,
where $r(n)=\lg\lg(n+4)/\sqrt{\lg(n+4)}$. That is, $m=o(n\sqrt{\lg n}/\lg\lg n)$ edges suffice. This includes $m\le n(\lg n)^{1/2-\theta}$ (\lean{hc_parity_sparse})
and every degree bound $\Delta$ (\lean{hc_parity_bounded}).
\end{theorem}

\begin{proof}[Proof sketch]
Algorithm~\ref{alg:hc} is Cut \& Count~\cite{CNPPRW11}. A chosen 2-factor with $c$ cycles has
$2^{c-1}$ side maps with $s_0=0$, which is odd iff $c=1$ (\lean{parity_sides}: flip the sides on a
cycle that avoids vertex 0), so $\#\psi\equiv\mathrm{hc}(G)\pmod2$ (\lean{hc_parity}). The sequential
counter has the same solutions on the edge and side bits (\lean{cntN_hcS}). For A5a$'$ the bound of
Theorem~\ref{thm:count} with width 3 and $N=5m+n+2$ variables is $c\,2^{9cN\,r(N)}$, with $r$ as in the theorem. Since $m\le n^2$, we have $N+4\le(n+4)^3$ and hence $r(N)\le3r(n)$
(\lean{rate_three}). So the exponent is at most $27c(5n\,g(n)+n+2)\,r(n)$, which is below
$\varepsilon n/2$ for large $n$ because $g(n)r(n)\to0$ and $r(n)\to0$ (\lean{boundO_le}).
\end{proof}

\begin{theorem}[\lean{ham_randomized}]\label{thm:hamrand}
There is one machine such that for every $D$, $\delta>0$ and $\varepsilon>0$ there are $k,t,n_0$
with the following property. For every graph $G$ with $n\ge n_0$ vertices and at most $Dn$ edges and
every seed $\omega$ of checks on the $m$ edge bits, the machine halts on $G$ and $\omega$ within
$2^{\varepsilon n}$ steps with a bit $\mathrm{out}(\omega)$. If $G$ has no Hamiltonian cycle,
$\mathrm{out}(\omega)=0$ for every seed; if it has one, $\mathrm{out}(\omega)=1$ for at least a
$2^{-\delta n}$ fraction of the seeds.
\end{theorem}

\begin{proof}[Proof sketch]
Add the checks and fixings of $\omega$ on the edge bits to the formula of A5a$'$, as in
Algorithm~\ref{alg:connrand} but without complement bits, since \COUNT{} does not fix the number of
true variables. The count is $\equiv$ the number of Hamiltonian cycles that pass the checks and
fixings (\lean{iso_parity}), and Theorem~\ref{thm:iso} applied to the set of Hamiltonian cycles in
$\{0,1\}^m$ leaves exactly one for a $2^{-\delta m}$ fraction of the seeds (\lean{iso_decide}). The
machine decodes $k$, the checks and the fixings from the input (\lean{decode_seedCode}), so one
machine serves every $\delta$. The formula has $5m+n+2+(3m+3)k=O(n)$ variables. Repetition over
$R$ seeds amplifies as in Section~\ref{sec:conn} (\lean{ham_amplified}).
\end{proof}

%======================================================================
\section{Parameterized by the solution size}\label{sec:fpt}

\begin{algorithm}[t]
\caption{Vertex Cover (A3). Lean: \lean{Vc.stepV}, \lean{Vc.leafV}, \lean{VcTM.vcM}.}\label{alg:vc}
\begin{algorithmic}[1]
\Require graph $G$, budget $k$; a degree $d$ chosen from $\varepsilon$
\State $\mathit{total}\gets0$
\For{each node $(k',E')$ of the tree from $(k,E(G))$, depth first}
  \If{$k'>0$ and an end $v$ of an edge of $E'$ has $\ge d$ neighbours $N(v)$}
    \State children: $(k'-1,\ E'\setminus\delta(v))$ and $(k'-|N(v)|,\ E'\setminus\delta(N[v]))$
  \Else
    \State $V'\gets$ the ends of the edges of $E'$
    \If{$|V'|\le k'd$} $\mathit{total}\gets\mathit{total}+\SIZED_2(|V'|,\min(k',|V'|),\{x_u\vee x_v:uv\in E'\})$
    \EndIf
  \EndIf
\EndFor
\State \Return $[\mathit{total}>0]$
\end{algorithmic}
\end{algorithm}

\begin{theorem}[\lean{vc_fpt}]
For every $\varepsilon>0$, one machine, given a graph $G$ on $n$ vertices (of any density) and $k\le
n$, decides whether $G$ has a vertex cover of at most $k$ vertices within $C(n+2)^C2^{\varepsilon k}$
steps.
\end{theorem}

\begin{proof}[Proof sketch]
Algorithm~\ref{alg:vc}; no kernel is used. A node has a cover within budget iff one child has
(\lean{good_step}), since a cover misses $v$ only if it contains $N(v)$. Along a path the budget drops
by 1 or by at least $d$, so by Anthropic's tree bound there are at most
$\sum_{j\le(k+d)/d}\binom{k+d}{j}\le2^{\varepsilon'(k+d)}$ leaves (\lean{leaves_le}). At a leaf all
degrees are below $d$, so a cover of size $\le k'$ leaves at most $k'd$ vertices on the edges
(\lean{card_V_le}), and $\SIZED_2$ decides the leaf in $2^{\delta kd}$ steps
(\lean{cnt_leafV_pos}); $\delta=\varepsilon/(4d)$ makes this $2^{\varepsilon k/4}$.
Lemma~\ref{lem:host} turns the walk into one machine.
\end{proof}

The best known algorithm for Vertex Cover runs in $O(1.2738^k+kn)$ time~\cite{CKX10}. The
Nonblocker problem asks for $k$ vertices whose complement dominates, i.e.\ for a dominating set of
at most $n-k$ vertices.

\begin{algorithm}[t]
\caption{Nonblocker (A3$'$). Lean: \lean{NbTM.preN}, \lean{NbTM.nbM}.}\label{alg:nb}
\begin{algorithmic}[1]
\Require graph $G$ on $n$ vertices, $k\le n$
\State $J\gets$ the vertices with a neighbour; $n'\gets|J|$
\If{$2k\le n'$} \Return yes \Comment{Ore: $G[J]$ has a dominating set of size $\le n'/2$}
\ElsIf{$n'<k$} \Return no
\Else\ \Return $[\DOMCOUNT(G[J],\ n'-k)>0]$ \Comment{$n'<2k$ vertices}
\EndIf
\end{algorithmic}
\end{algorithm}

\begin{theorem}[\lean{nb_fpt}]
For every $\varepsilon>0$, one machine, given $G$ on $n$ vertices and $k\le n$, decides whether $G$
has a dominating set of at most $n-k$ vertices within $C(n+2)^C2^{\varepsilon k}$ steps.
\end{theorem}

\begin{proof}[Proof sketch]
Isolated vertices lie in every dominating set, so the question is the same for $G[J]$ and the bound
$n'-k$ (\lean{nb_iff}). $G[J]$ has no isolated vertex, so by Ore's theorem~\cite{Ore62} it has a
dominating set of at most $n'/2$ vertices: the complement of a minimum dominating set dominates
(\lean{ore}). Supersets of dominating sets dominate, so in the remaining case a set of size exactly
$n'-k$ exists iff one of size at most $n'-k$ does (\lean{nb_answer}). \DOMCOUNT{} at $\delta=\varepsilon/2$
runs in $c\,2^{(\varepsilon/2)n'}\le c\,2^{\varepsilon k}$ steps.
\end{proof}

This is the linear kernel of~\cite{DFFPR06} in its simplest form, with the claimed refutation as the
exponential part.

%======================================================================
\section{Formalization and reference implementation}\label{sec:formal}

\paragraph{The Lean development.} The development consists of 45 Lean files with about 14{,}050
lines on top of Anthropic's development (pinned to a fixed revision) and Mathlib. A check file prints
the axioms of 78 theorems, the main theorems and the lemmas cited in the proof sketches; each depends on \lean{propext}, \lean{Classical.choice} and
\lean{Quot.sound} only, and there is no \lean{sorry}. Every main theorem quantifies over the
machine existentially and states its finite alphabets, its output (\lean{Wrap.Outputs}) and its
time bound explicitly. The polynomial-time front ends are verified with Anthropic's \lean{PolyT}
combinators and \lean{fun_prop}. Most of the proof effort goes into two kinds of lemmas: shape
bounds (width, number of variables and clauses of the formula, as functions of the input) and count
identities (the sized count equals the intended quantity), which are proved as explicit bijections
of finite sets.

\paragraph{Checks against Mathlib.} Lean cannot check that a definition means what is intended, so
where Mathlib has an independent definition we prove agreement. The number of Hamiltonian cycles
$\mathrm{hc}(G)$ is positive iff $G$ satisfies Mathlib's \lean{SimpleGraph.IsHamiltonian} (for
$n\ne1$), and it equals the number of edge sets of walks with Mathlib's \lean{IsHamiltonianCycle}
(\lean{numHamCycles_pos_iff}, \lean{numHamCycles_eq}); the proof uses Mathlib's result that a
connected graph with all degrees 2 is a single cycle. The homomorphism count equals the number of
Mathlib homomorphisms $G\to_g H$ for simple $H$ (\lean{numHom_simple}), both permanents are
stated with Mathlib's \lean{Matrix.permanent}, and the spanning trees of N7$''$ are subgraphs with
Mathlib's \lean{IsTree}. The remaining definitions (the spectrum, the Potts
counts, the Tutte polynomial, the seeds and the input codes) must be read by hand, as must
Anthropic's model of computation and statement files.

\paragraph{What is not formalized.} The numbered theorems, Lemma~\ref{lem:host}, the three
counters of Section~\ref{sec:prelim}, the special cases (independent sets, the Max-Cut spectrum, the
Ising model, the chromatic polynomial and $T(G;0,1-q)$) and the amplification of the one-sided error
are Lean theorems; the proof sketches summarize the Lean proofs. Not formalized are the results we
cite from the literature (lower bounds, earlier algorithms, Tutte's identity between $T(G;0,1-q)$ and
nowhere-zero flows) and the novelty labels of Table~\ref{tab:results}.

\paragraph{The reference implementation.} A Python package of about 2{,}150 lines mirrors every
front end and post-step clause by clause, with the same variable indices and from the same input
words as the Lean machines. Anthropic's counters are replaced by an exact DPLL counter with
component caching that returns the weight enumerator of the solutions, so the implementation runs
in exponential time; its purpose is to make the reductions concrete and testable. The tests
compare every algorithm with a direct enumeration of the quantity in the Lean statement (colorings,
homomorphisms, assignments, covers, dominating sets, permutations, Hamiltonian cycles, connected
vertex sets, connected dominating sets, spanning trees), check the shape bounds that the Lean proofs
use, and measure the one-sided error of A5b, N2$'$, N7$'$ and N7$''$ over random seeds. A further test checks that the Python front ends are the Lean ones:
a small Lean program evaluates the Lean definitions of all 20 front ends on 113 random input words,
and the formulas agree clause for clause.

\paragraph{Availability.} The Lean development, the Python package and the instructions to check
both are at \url{https://github.com/ssenge/eth-corollaries}.

%======================================================================
\section{Limitations and conclusion}\label{sec:concl}

\paragraph{What the theorems depend on.} Our results are Lean theorems; they do not assume an
unproved hypothesis. They are exactly as sound as Lean's kernel and Anthropic's formalization, and
they say what is intended only if Anthropic's statement files and the model of computation do. At
the time of writing nobody has reviewed these statements independently, and the claimed refutations
themselves are under review. The algorithms are galactic: no gain over $2^n$ is expected below
$n\approx10^{10}$.

\paragraph{What the results show.} The variable count is the currency. $O(n)$ variables give
$2^{\varepsilon n}$ for every $\varepsilon$; $n\cdot r(n)$ variables with $r=o(\sqrt{\lg n}/\lg\lg
n)$ still give $2^{o(n)}$; for the uniform bound the clause width must stay below about $(\lg
n)^{1/4}$. A sized counter makes the number of clauses irrelevant, which is why dense graphs are
possible for N2, N2$'$, N7, N7$'$, N8, N8$'$, N8$''$, A3 and A3$'$, and complement bits make it compatible with auxiliary
variables.

\paragraph{Open problems.} Colorings and homomorphisms into targets with $\lg h\ge(\lg
n)^{1/6}$; \#CNF-SAT of unbounded width with more clauses than in A1$'$; a deterministic version of
N2$'$, N7$'$ and N7$''$, that is, deterministic Steiner tree, connected vertex cover, connected
dominating set and Max Leaf Spanning Tree on dense graphs; the exact
number of connected vertex sets of size $\ell$ on dense graphs, of which N2 gives the parity;
Hamiltonian cycle on graphs of average degree $2^{\Omega(\sqrt{\lg n})}$, beyond Tulsyan's
bound~\cite{Tul26}; positive CNFs of unbounded width with superlinearly many clauses, such as Group
Steiner tree with $\omega(n)$ groups; the Tutte
polynomial off the integer hyperbolas; permanents with larger integer entries; the matching
polynomial of general graphs; and problems for which the refutations remove a lower bound but no
reduction with few variables is known, such as dense Max-Cut, CNF-SAT with $\poly(n)$ clauses,
planar problems below $2^{O(\sqrt n)}$ and Subset Sum.

\section*{Use of AI}

Claude Opus 5.5 (Anthropic) was used throughout this research: to explore approaches, to develop
and check proofs, to write the Lean formalization and the reference implementation, and to draft
the text. All theorems were checked by Lean.

\begin{thebibliography}{99}
\bibitem{anthropic} Anthropic. A $2^{o(n)}$ time algorithm for 3-SAT. Preprint hexagon:2610.00577v1,
  26 October 2026.
\bibitem{formalmath} Anthropic. formal-math: Lean formalizations, subdirectory \texttt{eth}.
  \url{https://github.com/anthropics/formal-math}, 2026.
\bibitem{Calabro09} C.~Calabro. \emph{The Exponential Complexity of Satisfiability Problems}.
  PhD thesis, UC San Diego, 2009.
\bibitem{CWY00} Y.~Caro, D.~B.~West, and R.~Yuster. Connected domination and spanning trees with
  many leaves. \emph{SIAM J.\ Discrete Math.}, 13(2):202--211, 2000.
\bibitem{CJ03} L.~Cai and D.~Juedes. On the existence of subexponential parameterized algorithms.
  \emph{J.\ Comput.\ System Sci.}, 67(4):789--807, 2003.
\bibitem{CIKP08} C.~Calabro, R.~Impagliazzo, V.~Kabanets, and R.~Paturi. The complexity of Unique
  $k$-SAT: an isolation lemma for $k$-CNFs. \emph{J.\ Comput.\ System Sci.}, 74(3):386--393, 2008.
\bibitem{CIP06} C.~Calabro, R.~Impagliazzo, and R.~Paturi. A duality between clause width and
  clause density for SAT. In \emph{CCC}, 2006.
\bibitem{CKX10} J.~Chen, I.~A.~Kanj, and G.~Xia. Improved upper bounds for vertex cover.
  \emph{Theoret.\ Comput.\ Sci.}, 411(40--42):3736--3756, 2010.
\bibitem{CFGKMPS17} M.~Cygan, F.~V.~Fomin, A.~Golovnev, A.~S.~Kulikov, I.~Mihajlin, J.~Pachocki,
  and A.~Soca{\l}a. Tight lower bounds on graph embedding problems. \emph{J.\ ACM}, 64(3):18, 2017.
\bibitem{CNPPRW11} M.~Cygan, J.~Nederlof, M.~Pilipczuk, M.~Pilipczuk, J.~M.~M.~van Rooij, and
  J.~O.~Wojtaszczyk. Solving connectivity problems parameterized by treewidth in single exponential
  time. In \emph{FOCS}, 2011.
\bibitem{dBKB26} M.~de Berg and S.~Kisfaludi-Bak. Lower bounds for dominating set in ball graphs and
  for weighted dominating set in unit-ball graphs. arXiv:2601.08425, 2026.
\bibitem{DFFPR06} F.~Dehne, M.~Fellows, H.~Fernau, E.~Prieto, and F.~Rosamond. Nonblocker:
  parameterized algorithmics for minimum dominating set. In \emph{SOFSEM}, LNCS~3831, 2006.
\bibitem{DHMTW14} H.~Dell, T.~Husfeldt, D.~Marx, N.~Taslaman, and M.~Wahl\'en. Exponential time
  complexity of the permanent and the Tutte polynomial. \emph{ACM Trans.\ Algorithms}, 10(4):21,
  2014.
\bibitem{FK72} C.~M.~Fortuin and P.~W.~Kasteleyn. On the random-cluster model I. Introduction and
  relation to other models. \emph{Physica}, 57(4):536--564, 1972.
\bibitem{GdHKS16} R.~Ganian, R.~de Haan, I.~Kanj, and S.~Szeider. On existential MSO and its
  relation to ETH. In \emph{MFCS}, LIPIcs 58:42, 2016.
\bibitem{IP01} R.~Impagliazzo and R.~Paturi. On the complexity of $k$-SAT. \emph{J.\ Comput.\
  System Sci.}, 62(2):367--375, 2001.
\bibitem{IPZ01} R.~Impagliazzo, R.~Paturi, and F.~Zane. Which problems have strongly exponential
  complexity? \emph{J.\ Comput.\ System Sci.}, 63(4):512--530, 2001.
\bibitem{JLSU14} P.~Jonsson, V.~Lagerkvist, J.~Schmidt, and H.~Uppman. Relating the time
  complexity of optimization problems in light of the exponential-time hypothesis. In
  \emph{MFCS}, LNCS 8635, pp.~408--419, 2014.
\bibitem{mathlib} The mathlib Community. The Lean mathematical library. In \emph{CPP}, 2020.
\bibitem{lean4} L.~de Moura and S.~Ullrich. The Lean 4 theorem prover and programming language. In
  \emph{CADE}, LNCS~12699, 2021.
\bibitem{openai} OpenAI. The Exponential Time Hypothesis is false. Preprint, 17 October 2026,
  revised 25 October 2026. \url{https://github.com/openai/math}, family 378.
\bibitem{Ore62} O.~Ore. \emph{Theory of Graphs}. AMS Colloquium Publications~38, 1962.
\bibitem{Ryser63} H.~J.~Ryser. \emph{Combinatorial Mathematics}. Carus Mathematical
  Monographs~14, Mathematical Association of America, 1963.
\bibitem{Tutte54} W.~T.~Tutte. A contribution to the theory of chromatic polynomials.
  \emph{Canad.\ J.\ Math.}, 6:80--91, 1954.
\bibitem{Tul26} V.~Tulsyan. Counting Hamiltonian cycles in sparse graphs in subexponential time,
  assuming the new \#SAT algorithms. Preprint and Lean development, October 2026; versions of
  28 and 29 October (commits ae323a3, ba4c161). \url{https://github.com/xennygrimmato/eth-hamiltonicity}.
\bibitem{VV86} L.~G.~Valiant and V.~V.~Vazirani. NP is as easy as detecting unique solutions.
  \emph{Theoret.\ Comput.\ Sci.}, 47:85--93, 1986.
\end{thebibliography}

\end{document}
